\par\Needspace{6\baselineskip}\medskip
\begin{center}\hypertarget{native-200-article150}{}\pdfbookmark[2]{150.}{article-150}150.\end{center}
Progredimur iam ad illustrationem huius methodi per aliquot exempla
copiose explicanda, quae simul evidentissime ostendent, quam late pateat, et
quam commode et expedite semper ad finem exoptatum perducat\footnote{Male loquuntur, qui methodum aliquam alia \textit{magis minusve exactam} pronunciant. Ea enim sola methodus problema solvisse censeri potest, per quam quemvis praecisionis gradum attingere saltem in potestate est. Quamobrem methodus alia alii eo tantum nomine palmam praeripit, quod \textit{eundem} praecisionis gradum per aliam celerius minorique labore, per aliam tardius graviorique opera assequi licet.}.

Exemplum \textit{primum} planeta novus Iuno nobis suppeditabit, ad quem finem
observationes sequentes Grenovici factas et a cel. \textsc{Maskelyne} nobiscum communicatas
eligimus.
\SourceLeaf{201}
\begin{center}
\small\setlength{\tabcolsep}{5pt}\renewcommand{\arraystretch}{1.2}
\begin{tabular}{|r|c|c|}
\hline
Temp. med. Grenov. & Ascens. recta app. & Decl. austr. app.\\
\hline
1804 Oct. 5.\ $10^{\mathrm{h}}\,51^{\mathrm{m}}\,6^{\mathrm{s}}$ & $357^{0}\,10'\,22''{,}35$ & $6^{0}\,40'\,8''$\\
17.\quad $9\quad58\quad10$ & $355\quad43\quad45{,}30$ & $8\quad47\quad25$\\
27.\quad $9\quad16\quad41$ & $355\quad11\quad10{,}95$ & $10\quad2\quad28$\\
\hline
\end{tabular}
\end{center}

E tabulis Solaribus pro iisdem temporibus invenitur
\begin{center}
\footnotesize\setlength{\tabcolsep}{3pt}\renewcommand{\arraystretch}{1.2}
\begin{tabular}{|r|c|c|c|c|c|}
\hline
 & \shortstack{Longit. Solis ab\\aequin. appar.} & Nutatio & \shortstack{Distantia a\\terra} & \shortstack{Latitudo\\Solis} & \shortstack{Obliquitas appar.\\eclipticae}\\
\hline
Oct. 5 & $192^{0}\,28'\,53''{,}71$ & $+15''{,}43$ & $0{,}998\,8839$ & $-0''{,}49$ & $23^{0}\,27'\,59''{,}48$\\
17 & $204\quad20\quad21{,}54$ & $+15{,}51$ & $0{,}995\,3968$ & $+0{,}79$ & $59{,}26$\\
27 & $214\quad16\quad52{,}21$ & $+15{,}60$ & $0{,}992\,8340$ & $-0{,}15$ & $59{,}06$\\
\hline
\end{tabular}
\end{center}

Calculum ita adstruemus, ac si orbita adhuc penitus incognita esset:
quamobrem loca Iunonis a parallaxi liberare non licebit, sed hanc ad loca
terrae transferre oportebit. Primo itaque ipsa loca observata ab aequatore ad
eclipticam reducimus, adhibita obliquitate apparente, unde prodit:
\begin{center}
\small\setlength{\tabcolsep}{6pt}\renewcommand{\arraystretch}{1.2}
\begin{tabular}{|r|c|c|}
\hline
 & \shortstack{Longit. appar.\\Iunonis} & \shortstack{Latit. appar.\\Iunonis}\\
\hline
Oct. 5 & $354^{0}\,44'\,54''{,}28$ & $-4^{0}\,59'\,31''{,}59$\\
17 & $352\quad34\quad44{,}51$ & $-6\quad21\quad56{,}25$\\
27 & $351\quad34\quad51{,}57$ & $-7\quad17\quad52{,}70$\\
\hline
\end{tabular}
\end{center}

Cum hoc calculo statim iungimus determinationem longitudinis et latitudinis
ipsius zenith loci observationis in tribus observationibus: rectascensio quidem
cum rectascensione Iunonis convenit (quod observationes in ipso meridiano
sunt factae), declinatio autem aequalis est altitudini poli $=51^{0}\,28'\,39''$.
Ita obtinemus
\begin{center}
\small\setlength{\tabcolsep}{6pt}\renewcommand{\arraystretch}{1.2}
\begin{tabular}{|r|c|c|}
\hline
 & Long. ipsius zenith & Latitudo\\
\hline
Oct. 5 & $24^{0}\,29'$ & $46^{0}\,53'$\\
17 & $23\quad25$ & $47\quad24$\\
27 & $23\quad1$ & $47\quad36$\\
\hline
\end{tabular}
\end{center}
\noindent\SourceLeaf{202}Iam ad normam praeceptorum in art. 72 traditorum determinabuntur terrae
loci ficti in ipso plano eclipticae, in quibus corpus coeleste perinde apparuisset
atque in locis veris observationum. Hoc modo prodit, statuendo parallaxin
Solis mediam $=8''{,}6$
\begin{center}
\small\setlength{\tabcolsep}{4pt}\renewcommand{\arraystretch}{1.2}
\begin{tabular}{|r|c|c|c|}
\hline
 & Reductio longit. & Reductio distantiae & Reductio temporis\\
\hline
Oct. 5 & $-22''{,}39$ & $+0{,}000\,3856$ & $-0^{\mathrm{s}}{,}19$\\
17 & $-27{,}21$ & $+0{,}000\,2329$ & $-0{,}12$\\
27 & $-35{,}82$ & $+0{,}000\,2085$ & $-0{,}12$\\
\hline
\end{tabular}
\end{center}
\noindent Reductio temporis ideo tantum adiecta est, ut appareat, eam omnino insensibilem
esse.

Deinde omnes longitudines tum planetae tum terrae reducendae sunt ad
aequinoctium vernale medium pro aliqua epocha, pro qua adoptabimus initium
anni 1805; subducta itaque nutatione adhuc adiicienda est praecessio, quae
pro tribus observationibus resp. est $11''{,}87$, $10''{,}23$, $8''{,}86$, ita ut pro observatione
prima addere oporteat $-3''{,}56$, pro secunda $-5''{,}28$, pro tertia $-6''{,}74$.

Denique longitudines et latitudines Iunonis ab aberratione fixarum purgandae
sunt; sic per regulas notas invenitur, a longitudinibus resp. subtrahi
debere $19''{,}12$, $17''{,}11$, $14''{,}82$, latitudinibus vero addi $0''{,}53$, $1''{,}18$, $1''{,}75$, per quam
additionem valores absoluti diminutionem patientur, quoniam latitudines australes
tamquam negativae spectantur.
